% 20260522_Assingment_Probability_Statistics
\documentclass[dvipdfmx]{jsreport}
\usepackage{soupreamble}

\renewcommand{\P}{\mathcal{P}}
%\renewcommand{\cB}{\mathcal{B}}

\title{確率及び統計学特論A 第2回課題}
\author{M1 6012 川村聡一郎}
\date{}

\begin{document}
\maketitle

\begin{tcolorbox}
	確率変数$X, Y$に対して
\begin{equation*}
	\mathcal{P}(X\le a)= \mathcal{P}(Y\le a),\ a\in \R
\end{equation*}
	ならば，
\begin{equation*}
	\mathcal{P}(X\in A)= \mathcal{P}(Y\in A),\ A\in \mathcal{B}(\R)
\end{equation*}
	であることを示せ．
\end{tcolorbox}


\begin{souanswer} %答案はここに書く
	実数$a\in \R$に対して，区間$I_a= (-\infty, a]$を考える．このような区間全体からなる集合族を$\P$とおく．
\begin{equation}
	\P= \{(-\infty, a]\; a\in \R\}
\end{equation}
	$a, b\in\R$に対して$(-\infty, a]\cap (-\infty, b]= (-\infty, \min{(a, b)}]\in \P$であるためには$\P$は共通部分について閉じている．
	すなわち$\P$は$\pi$系であり，Borel集合族の定義より
\begin{equation}
	\sigma(\P)= \mathcal{B}(\R)
\end{equation}
	確率が一致するようなBorel集合の集合を$mathcal{L}$とする．
\begin{equation}
	\mathcal{L}= \{A\in \mathcal{B(\R)}; \P(X\in A)= \P(Y\in A)\}
\end{equation}
	この$\mathcal{L}$が$\lambda$系の条件を満たすことを示す．
\begin{enumerate_alpha}
	\item $\R\in \mathcal{L}$：
	$\P(X\in A)= 1$かつ$\P(Y\in A)= 1$であるから，
\begin{equation}
	\P(X\in A)= \P(Y\in A)
\end{equation}
	となり，$\R\in \mathcal{L}$が成り立つ．

	\item $A, B\in \mathcal{L}, A\subset B\Rightarrow B\setminus A\in \mathcal{L}$：
	$A, B\in \mathcal{L}, A\subset B$とする．確率の加法性より
\begin{align}
	\P(X\in B\setminus A)&= \P(X\in B)- \P(X\in A)\\
	\P(Y\in B\setminus A)&= \P(Y\in B)- \P(Y\in A)
\end{align}
	$A, B\in \mathcal{L}$の定義より，$\P(X\in A)= \P(Y\in A)$かつ$\P(Y\in B)= \P(Y\in B)$であるから
\begin{equation}
	\P(X\in B\setminus A)= \P(Y\in B\setminus A)
\end{equation}
	であり，$B\setminus A\in \mathcal{L}$
	\item $\{A_n\}\in \mathcal{L}$が単調増加ならば$\bigcup_{n= 1}^\infty A_n\in \mathcal{L}$：
	互いに素な可算個のBorel集合列$\{A_n\}\subset \mathcal{L}$をとる．確率の可算加法性より
\begin{align}
	\P\left(X\in \bigcup_{n= 1}^\infty A_n\right)&= \sum_{n= 1}^\infty \P(X\in A_n)\\
	\P\left(Y\in \bigcup_{n= 1}^\infty A_n\right)&= \sum_{n= 1}^\infty \P(Y\in A_n)
\end{align}
	各$n$について$A_n\in \mathcal{L}$であり，$\P(X\in A_n)= \P(Y\in A_n)$であるから各項がすべて等しく，無限級数の和も等しくなる．
\begin{equation}
	\P\left(X\in \bigcup_{n= 1}^\infty A_n\right)= \P\left(Y\in \bigcup_{n= 1}^\infty A_n\right)
\end{equation}
	となり，$\bigcup_{n= 1}^\infty A_n\in \mathcal{L}$
	よって$\mathcal{L}$は$\lambda$系
\end{enumerate_alpha}
	任意の$a\in \R$に対して$\P(X\le a)= \P(Y\le a)$が成り立つ．これは$\P\subset \mathcal{L}$を意味する．ここで$\P$は$\pi$系，$\mathcal{L}$は$\lambda$系である．$\pi\mathrm{-}\lambda$定理より
\begin{equation}
	\sigma(\P)\subset \mathcal{L}
\end{equation}
	が成り立ち，$\sigma(\P)= \mathcal{B}(\R)$であるから
\begin{equation}
	\mathcal{B}(\R)\subset \mathcal{L}
\end{equation}
	が得られ，$\mathcal{L}$の定義より任意の$A\in \mathcal{B}(\R)$に対して$\P(X\le a)= \P(Y\le a)$が成立．
\end{souanswer}
\end{document}
